% Section 1 -- Introduction
% Target: 2 pages.
% Goal: motivate the problem, state the contribution claim, outline the paper.

The agent economy of 2026 has a delegation problem.  When an AI agent
acts on behalf of a human, who then delegates to another agent, who
delegates to a third, the question \emph{who actually authorized
what?} becomes load-bearing for security, legal accountability, and
auditing.  The standards layer has responded: OAuth 2.0 Token
Exchange (RFC~8693)~\cite{rfc8693} provides nested \texttt{act}
claims; the IETF~\texttt{draft-ietf-oauth-attenuating-agent-tokens}
formalizes attenuation; W3C Verifiable Credentials provides the
data-model substrate.  The research layer has responded too:
Authenticated Delegation (South et al., ICML
2025)~\cite{south2025authdelegation}; AITH (Chen,
2026)~\cite{chen2026aith} mechanizes five Tamarin theorems for
post-quantum delegation; LLMbda Calculus (Garby et al.,
2026)~\cite{garby2026llmbda} proves termination-insensitive
noninterference for an LLM-invoking lambda calculus.

Yet a vulnerability surfaced in March 2026 on the IETF OAuth mailing
list~\cite{ietf2026splicing}, \emph{delegation chain splicing}, remains
imperfectly addressed.  An attacker presenting mismatched
\texttt{subject\_token} and \texttt{actor\_token} inputs to an STS
that validates each independently can extract a token asserting a
delegation chain that never occurred.  The structural fix---requiring
the chain's middle principal to appear in both halves---is well
understood, but no existing system enforces it through a unified
mechanism that spans the logic, the protocol, and the wire format.

This paper introduces \emph{DLC} (Delegation Logic Calculus), a
formal system in which the same-principal condition is enforced by
construction across four independent axes:

\begin{itemize}
\item \textbf{In the type theory}: the \textsc{delegate} typing rule
  has a single $q$ variable bound in both premises (Section~\ref{sec:calculus}).
\item \textbf{In the symbolic protocol model}: the
  \texttt{Delegate\_Accept} multiset-rewriting rule's input pattern
  unifies the principal across both signature consumptions
  (Section~\ref{sec:symbolic}).
\item \textbf{In the cryptographic correspondence (T2)}: an EUF-CMA
  adversary spanning a spliced chain reduces to forgery of one of the
  two underlying signatures, which by Ed25519's EUF-CMA security is
  negligible (Section~\ref{sec:computational}).
\item \textbf{In the wire format}: each delegation step is content-
  addressed by SHA-256 of its canonical bytes, so any chain-
  splicing attempt produces detectable hash discontinuities
  (Section~\ref{sec:wire}).
\end{itemize}

The contribution is the \emph{diagonal}: no existing system spans all
four axes with mechanical verification.  AITH covers the symbolic
axis only; LLMbda covers the type-theoretic axis only; macaroons
and biscuits cover the wire axis only; OAuth drafts cover the
standards axis only.  DLC's unification, with the chain-splicing-
defeating rule appearing in identical form at every layer, is what
makes it a candidate for the role of a \emph{citable foundation}
for delegation work in the agent economy.

\paragraph{Contributions}

This paper makes the following claims:

\begin{enumerate}
\item A modal-linear-temporal authorization logic with a no-chain-
  splicing delegation rule, mechanized in Lean~4 with no
  \texttt{sorry} on the proofs that are closed (T4 non-introduction,
  graded comonad laws, indexed monad laws, four structural
  substitution lemmas).
\item A symbolic protocol model verified independently in Tamarin
  and ProVerif, demonstrating that the unification of two
  symbolic provers detects modeling gaps that single-prover work
  misses (Section~\ref{sec:symbolic}).
\item A Rust reference implementation with 76 unit tests, including
  a CBOR wire format with round-trip-tested encoding for all
  21 term constructors and Tamarin/ProVerif fact exporters.
\item A reproducible verification ledger (\texttt{make ledger})
  emitting a single JSON document containing the status of every
  theorem, every model, every drift check, and every test count.
\item An honest assessment (Section~\ref{sec:conclusion}) of what is
  proven, what is stated, and what remains for collaborator
  engagement.  We make no claim to having closed T2 unconditionally
  or T3 in full; both are stated and partially proven, with
  documented paths to closure via VCVio~\cite{tuma2026vcvio} for T2
  and hand-rolled logical relations for T3.
\end{enumerate}

\paragraph{Paper structure}

Section~\ref{sec:calculus} defines DLC's syntax and typing rules.
Section~\ref{sec:metatheory} states the four headline theorems and
reports the current proof status of each.  Section~\ref{sec:symbolic}
presents the Tamarin and ProVerif models.  Section~\ref{sec:wire}
specifies the wire format.  Section~\ref{sec:computational} sketches
the EUF-CMA reduction via VCVio.  Section~\ref{sec:related} situates
DLC against AITH, LLMbda, NAL, Garg-Pfenning, macaroons, and the
RFC~8693 lineage.  Section~\ref{sec:conclusion} discusses limitations.

\paragraph{Artifact}

The complete DLC artifact---Lean development, Rust implementation,
Tamarin and ProVerif models, EasyCrypt skeleton, IETF draft, W3C VC
context---is available at
\url{https://github.com/coproduct-opensource/delegation_calc}.  A
reproducible build under Nix produces a byte-identical
\texttt{ledger.json}; the same document is published as a release
artifact on every tagged version.
